% Options for packages loaded elsewhere
\PassOptionsToPackage{unicode}{hyperref}
\PassOptionsToPackage{hyphens}{url}
\documentclass[
  ignorenonframetext,
]{beamer}
\newif\ifbibliography
\usepackage{pgfpages}
\setbeamertemplate{caption}[numbered]
\setbeamertemplate{caption label separator}{: }
\setbeamercolor{caption name}{fg=normal text.fg}
\beamertemplatenavigationsymbolsempty
% remove section numbering
\setbeamertemplate{part page}{
  \centering
  \begin{beamercolorbox}[sep=16pt,center]{part title}
    \usebeamerfont{part title}\insertpart\par
  \end{beamercolorbox}
}
\setbeamertemplate{section page}{
  \centering
  \begin{beamercolorbox}[sep=12pt,center]{section title}
    \usebeamerfont{section title}\insertsection\par
  \end{beamercolorbox}
}
\setbeamertemplate{subsection page}{
  \centering
  \begin{beamercolorbox}[sep=8pt,center]{subsection title}
    \usebeamerfont{subsection title}\insertsubsection\par
  \end{beamercolorbox}
}
% Prevent slide breaks in the middle of a paragraph
\widowpenalties 1 10000
\raggedbottom
\AtBeginPart{
  \frame{\partpage}
}
\AtBeginSection{
  \ifbibliography
  \else
    \frame{\sectionpage}
  \fi
}
\AtBeginSubsection{
  \frame{\subsectionpage}
}
\usepackage{iftex}
\ifPDFTeX
  \usepackage[T1]{fontenc}
  \usepackage[utf8]{inputenc}
  \usepackage{textcomp} % provide euro and other symbols
\else % if luatex or xetex
  \usepackage{unicode-math} % this also loads fontspec
  \defaultfontfeatures{Scale=MatchLowercase}
  \defaultfontfeatures[\rmfamily]{Ligatures=TeX,Scale=1}
\fi
\usepackage{lmodern}
\ifPDFTeX\else
  % xetex/luatex font selection
\fi
% Use upquote if available, for straight quotes in verbatim environments
\IfFileExists{upquote.sty}{\usepackage{upquote}}{}
\IfFileExists{microtype.sty}{% use microtype if available
  \usepackage[]{microtype}
  \UseMicrotypeSet[protrusion]{basicmath} % disable protrusion for tt fonts
}{}
\makeatletter
\@ifundefined{KOMAClassName}{% if non-KOMA class
  \IfFileExists{parskip.sty}{%
    \usepackage{parskip}
  }{% else
    \setlength{\parindent}{0pt}
    \setlength{\parskip}{6pt plus 2pt minus 1pt}}
}{% if KOMA class
  \KOMAoptions{parskip=half}}
\makeatother
\usepackage{longtable,booktabs,array}
\usepackage{caption}
\captionsetup[table]{skip=6pt}
\newcounter{none} % for unnumbered tables
\usepackage{calc} % for calculating minipage widths
\usepackage{caption}
% Make caption package work with longtable
\makeatletter
\def\fnum@table{\tablename~\thetable}
\makeatother
\setlength{\emergencystretch}{3em} % prevent overfull lines
\providecommand{\tightlist}{%
  \setlength{\itemsep}{0pt}\setlength{\parskip}{0pt}}
% Auto-generated by infrastructure/rendering/_pdf_latex_helpers.py::extract_math_font_preamble.
% Minimal header passed to Pandoc via -H so Beamer slides pick up the same math font
% as the combined-PDF path. Adjust the manuscript's preamble.md to change the font globally.
\usepackage{unicode-math}
\setmathfont{latinmodern-math.otf}

% Manuscript macro/package fallbacks (newcommand -> providecommand).
% Core mathematics
\usepackage{amsmath}
\usepackage{amssymb}
% Document layout
% Tables
\usepackage{booktabs}
% Algorithm typesetting (for pseudocode in §2 Methodology)
% Code listings
\usepackage{listings}
% Typography and formatting
% Cross-references and citations
% ── Unicode-capable mono font for code listings ──────────────────────
% JuliaMono (TeX Live 2026) covers the full math/Greek glyph set used in
% scientific code blocks; the default lmmono lacks \alpha/\mu/\partial/\nabla.
%
% The hardcoded TeX Live 2026 path below is portable to most macOS / Linux
% TeX Live installs that placed JuliaMono under the default texmf-dist path.
% If your TeX Live version differs (e.g., 2025, 2027) or your distribution
% installs fonts elsewhere, comment out the explicit Path = line — fontspec
% will then search the system font path. If JuliaMono is not installed at
% all, comment out the whole block; xelatex falls back to lmmono with
% reduced Unicode coverage but the document still compiles.
% Math font for unicode-math: Latin Modern Math (TeX Live) has full BMP
% coverage including U+2223 (\mid), U+226A/226B (\ll/\gg), and the Greek/
% blackboard letters used in equations. Without an explicit \setmathfont,
% unicode-math falls back to lmroman text font which lacks several glyphs
% and emits "Missing character" warnings on every \mid in math mode.

% Slide layout defaults for warning-clean scientific decks.
\usepackage{etoolbox}
\IfFileExists{xurl.sty}{\usepackage{xurl}}{}
\IfFileExists{seqsplit.sty}{\usepackage{seqsplit}}{\newcommand{\seqsplit}[1]{#1}}
\protected\def\breakseq#1{\seqsplit{#1}}
\protected\def\breaktt#1{\begingroup\ttfamily\seqsplit{#1}\endgroup}
\setlength{\emergencystretch}{6em}
\tolerance=5000
\hbadness=10000
\hfuzz=1pt
\setlength{\tabcolsep}{2pt}
\AtBeginEnvironment{longtable}{\tiny\renewcommand{\arraystretch}{0.86}\setlength{\tabcolsep}{1pt}}
\AtBeginEnvironment{tabular}{\tiny\renewcommand{\arraystretch}{0.86}\setlength{\tabcolsep}{1pt}}
\AtBeginEnvironment{equation}{\tiny}
\AtBeginEnvironment{equation*}{\tiny}
\AtBeginEnvironment{align}{\tiny}
\AtBeginEnvironment{align*}{\tiny}
\AtBeginEnvironment{itemize}{\footnotesize}
\AtBeginEnvironment{enumerate}{\footnotesize}
\AtBeginEnvironment{description}{\footnotesize}
\setbeamerfont{caption}{size=\tiny}
\setbeamerfont{caption name}{size=\tiny}
\setbeamerfont{normal text}{size=\small}
\setbeamerfont{frametitle}{size=\small}
\setbeamerfont{section title}{size=\footnotesize}
\setbeamerfont{subsection title}{size=\footnotesize}
\setbeamertemplate{section page}{%
  \centering
  \begin{beamercolorbox}[sep=12pt,center,wd=\paperwidth]{section title}
    \parbox{0.86\paperwidth}{\centering\usebeamerfont{section title}\insertsection\par}
  \end{beamercolorbox}
}
\setbeamertemplate{subsection page}{%
  \centering
  \begin{beamercolorbox}[sep=8pt,center,wd=\paperwidth]{subsection title}
    \parbox{0.86\paperwidth}{\centering\usebeamerfont{subsection title}\insertsubsection\par}
  \end{beamercolorbox}
}
\setlength{\abovecaptionskip}{2pt}
\setlength{\belowcaptionskip}{0pt}

% Natbib and cross-reference fallbacks — slides don't load natbib
% or cleveref, but combined-PDF manuscript prose may emit these
% commands. The fallback renders citations as a bracketed key list
% and cross-references as detokenized labels so slides don't fail on
% undefined control sequences or raw underscores. \providecommand is
% a no-op if packages load later.
\providecommand{\citep}[1]{[#1]}
\providecommand{\citet}[1]{#1}
\providecommand{\citealp}[1]{#1}
\providecommand{\citeauthor}[1]{#1}
\providecommand{\citeyear}[1]{#1}
\providecommand{\cref}[1]{\texttt{\detokenize{#1}}}
\providecommand{\Cref}[1]{\texttt{\detokenize{#1}}}
\providecommand{\crefrange}[2]{\texttt{\detokenize{#1}}--\texttt{\detokenize{#2}}}
\providecommand{\Crefrange}[2]{\texttt{\detokenize{#1}}--\texttt{\detokenize{#2}}}
\providecommand{\paragraph}[1]{\textbf{#1}\ }

\newtheorem{proposition}{Proposition}
\newtheorem{hypothesis}{Hypothesis}
\newtheorem{remark}[theorem]{Remark}
\newtheorem{axiom}{Axiom}
\newtheorem{property}{Property}
\usepackage{bookmark}
\IfFileExists{xurl.sty}{\usepackage{xurl}}{} % add URL line breaks if available
\urlstyle{same}
% fallback for those not using the hyperref driver hyperxmp:
\makeatletter
\@ifundefined{xmpquote}{\newcommand{\xmpquote}[1]{#1}}{}
\makeatother
\hypersetup{
  hidelinks,
  pdfcreator={LaTeX via pandoc}}

\author{\texorpdfstring{}{}}
\date{}

\begin{document}

\section{Conventional Desktops: Hardening Candidates, Compatibility
Costs, and Agent Isolation}\label{sec:desktops}

\begin{frame}[fragile,allowframebreaks]{Conventional Desktops: Hardening
Candidates, Compatibility Costs, and Agent Isolation}
The two preceding sections examined the architectural poles of this
review: hypervisor-based compartmentalization in sec.~5
and declarative, rebuildable operations in sec.~6. Most
operators, however, will run a conventional Linux desktop for the
foreseeable future, so the desktop candidates deserve the same
property-based scrutiny rather than a distribution-label comparison.
This section evaluates seven desktop-focused candidates against the nine
properties of sec.~4, with particular
attention to \breaktt{application\_confinement}, \texttt{integrity}, and
\texttt{human\_usability}: whether shipped defaults actually constrain
applications, whether boot and update paths verify what runs, and
whether a real operator will sustain the configuration. All judgments
follow from documented designs, release notes, and advisories reviewed
on 2026-09-10; none rests on a comparative penetration test, and none is
\end{frame}

\begin{frame}[fragile,allowframebreaks]{Conventional Desktops: Hardening
Candidates, Compatibility Costs, and Agent Isolation}
expressed as a numeric score.

The desktop lane is also where the fastest documented movement in this
review is happening. The secureblue release sequence from v4.3.0 through
v4.9.1, Fedora's removal of the GNOME X11 session, Debian's LTS
planning, and Ubuntu's AppArmor and snap-permission changes all landed
inside the review window, and each moves a candidate's documented
defaults in a direction this section can now assess concretely rather
than in principle (secureblue Project 2026b, 2026c; Fedora Project
2026e; Project 2025; Canonical 2026).

\end{frame}

\begin{frame}[fragile,allowframebreaks]{Conventional Desktops: Hardening
Candidates, Compatibility Costs, and Agent Isolation}
\begin{longtable}[]{@{}
  >{\raggedright\arraybackslash}p{(\linewidth - 6\tabcolsep) * \real{0.2500}}
  >{\raggedright\arraybackslash}p{(\linewidth - 6\tabcolsep) * \real{0.2500}}
  >{\raggedright\arraybackslash}p{(\linewidth - 6\tabcolsep) * \real{0.2500}}
  >{\raggedright\arraybackslash}p{(\linewidth - 6\tabcolsep) * \real{0.2500}}@{}}
\caption{Desktop candidates compared by documented security-relevant
design, important limitation, and analytical assessment. Judgments are
analytical, derived from documented designs and project documentation;
they are not results of a comparative penetration
test.}\label{tbl:desktops}\tabularnewline
\toprule\noalign{}
\begin{minipage}[b]{\linewidth}\raggedright
Candidate
\end{minipage} & \begin{minipage}[b]{\linewidth}\raggedright
Documented security-relevant design
\end{minipage} & \begin{minipage}[b]{\linewidth}\raggedright
Important limitation
\end{minipage} & \begin{minipage}[b]{\linewidth}\raggedright
Assessment
\end{minipage} \\
\midrule\noalign{}
\endfirsthead
\toprule\noalign{}
\begin{minipage}[b]{\linewidth}\raggedright
Candidate
\end{minipage} & \begin{minipage}[b]{\linewidth}\raggedright
Documented security-relevant design
\end{minipage} & \begin{minipage}[b]{\linewidth}\raggedright
Important limitation
\end{minipage} & \begin{minipage}[b]{\linewidth}\raggedright
Assessment
\end{minipage} \\
\midrule\noalign{}
\endhead
\textbf{Fedora Silverblue / Kinoite} & Silverblue documents read-only
system paths with writable \texttt{/etc} and \texttt{/var}, and Fedora
enforces SELinux as its default baseline (Fedora Project 2026b, 2026c).
Fedora is also removing the GNOME X11 session: the Wayland-only GNOME
change strips X11 session packages from the Fedora 43 and 44 GNOME
offerings (Fedora Project 2026e). & Persistent writable state remains,
and approximately 13 months of release support require regular upgrades
to a new release, with Fedora 42 reaching end of life on 2026-05-13 and
Fedora 43 on 2026-12-09 (Fedora Project 2026a). Silverblue-specific
details should not be transplanted onto Kinoite without checking each
image. & A credible atomic-desktop direction inside the Fedora
ecosystem: image-based system updates, an enforcing MAC baseline by
default, and a display stack being reduced to one code path. Its
documented claims concern system paths and update mechanics, not
containment of a compromised application's access to that application's
own data. \\
\textbf{Fedora Workstation} & Fedora's documented baseline includes
SELinux enforcement without choosing an atomic variant, and the same
Wayland-only GNOME trajectory applies to the non-atomic desktop (Fedora
Project 2026c, 2026e). & SELinux policy and application sandboxing must
be evaluated as deployed; the existence of a kernel access-control
mechanism does not by itself establish a tight sandbox for each
application (Fedora Project 2026d). The X11 removal also narrows
fallback options for specialized workflows that still require X11. & A
defensible compatibility-first choice when the operator preserves the
baseline and places dangerous work in separate isolation. An atomic
variant does not automatically win every runtime-security comparison
against it. \\
\textbf{secureblue} & The project documents a Fedora Atomic base, broad
\texttt{hardened\_malloc} deployment, a confined Trivalent browser,
removal of SUID-root programs, restrictive application settings, and
signed-container policy (secureblue Project 2026a). Since v4.3.0, an
SELinux policy denies user-namespace creation to unconfined processes
and the container domain by default --- a default-deny posture with
explicit \texttt{ujust} toggles for Flatpak, Chromium, and container
workflows (secureblue Project 2026b). The v4.9.1 release completed a
SUID-less userland on 2026-05-26, removing the legacy setuid privilege
path structurally rather than by policy (secureblue Project 2026c). &
The same documentation records compatibility-affecting restrictions:
disabled Xwayland by default, default-deny user namespaces that require
per-domain toggles, and image renames that follow Fedora's own (the
browser rebranded to Trivalent in the v4.4.0 cycle) (secureblue Project
2026a). These are project-documented measures, not measured superiority
results. & The most compelling conventional security-focused desktop
candidate in this review, conditional on application compatibility,
update reliability, and the operator not undoing the hardening to
restore convenience. \\
\textbf{Debian stable} & Debian operates a coordinated security-advisory
process and recommends unattended security upgrades; LTS extends stable
support under a separate group from the Debian security team (D. Project
2026b, 2026a). Debian 13 ``trixie'' released 2025-08-09 with full
support into 2028 and LTS coverage to 2030-06-30; the Debian 14
``forky'' cycle plans to require reproducible package builds as a matter
of policy (Project 2025). & A maintenance policy is not evidence of
strong isolation between arbitrary desktop applications, and the two
support phases are organized differently; conflating them misstates who
maintains what. The reproducible-packages plan concerns build
provenance, not runtime confinement. & A sound conservative platform
when deliberately configured for the actual workload, now with the
longest explicitly dated LTS horizon in this table. A well-maintained,
narrowly configured installation is preferable to a more elaborate stack
the owner cannot sustain. \\
\textbf{Ubuntu LTS} & Ubuntu documents five years of standard security
maintenance for LTS releases, with longer coverage under specified
offerings; snap confinement distinguishes strict, classic, and
development modes (Ubuntu 2026b, 2026a). Ubuntu 26.04 LTS extends
AppArmor enforcement to more daemons (OpenLDAP among them) and ships an
experimental permission-prompting mechanism for strictly confined snaps
(Canonical 2026). & Classic snaps carry no confinement, so ``installed
as a snap'' is not sufficient evidence of application isolation;
coverage must be matched to the package, release, and support
entitlement. The permission-prompting experiment asks the user to
adjudicate requests --- a cognitive surface, not a kernel boundary ---
and its defaults still require review. & A pragmatic supported baseline
for workstation work, provided confinement mode and support entitlement
are checked per application rather than assumed from the packaging name.
The 26.04 direction moves Ubuntu toward mediated permissions, with the
mediation burden landing on the operator. \\
\textbf{Kicksecure} & Kicksecure documents Debian-based hardening and a
user/system-maintenance split with different boot roles; the split's
documentation labels verified boot as planned rather than shipped, with
a conditional design for validating the boot chain before granting it
trust (K. Project 2026a, 2026b, 2026c). & The hardening wiki explicitly
contains research material, non-default proposals, and changes that can
cause breakage; a long wiki page must not be mistaken for a list of
shipped defaults. Verified boot remains a roadmap item, not an installed
property. & A serious candidate where reducing everyday administrative
authority is valuable. Judge the installed release and enabled settings;
the split is an interesting direction for separating daily use from
system maintenance, with boot integrity still on the roadmap rather than
in the release. \\
\textbf{openSUSE Aeon} & Aeon describes an immutable desktop based on
Tumbleweed, while transactional-update prepares changes in a new
snapshot without modifying the running system; the project still labels
its releases release-candidate ({openSUSE} 2026a, 2026b). &
Snapshot-based system replacement is not evidence that a compromised
application cannot read or modify its authorized user data; the
documented mechanism principally concerns system updates. The
release-candidate label is itself a caveat an operator should not read
past. & A credible image/snapshot-oriented desktop direction that has
not completed its own stability review. This review did not verify
role-specific defaults for the closely related MicroOS server image, and
records those fields as n.a. rather than inferring them from Aeon's
design. \\
\bottomrule\noalign{}
\end{longtable}

\end{frame}

\begin{frame}[fragile,allowframebreaks]{Conventional Desktops: Hardening
Candidates, Compatibility Costs, and Agent Isolation}
Two honesty notes apply to the table. First, the Aeon entry deliberately
withholds claims about the openSUSE MicroOS server image: this review
did not verify all current MicroOS role-specific defaults or a distinct
boot-integrity baseline, and the analysis treats those fields as n.a.
rather than inferring server behavior from a desktop sibling that shares
the transactional-update mechanism ({openSUSE} 2026b). Second,
``immutable'' is treated throughout as a statement about particular
system paths and update mechanisms, never as a guarantee against runtime
data theft from a compromised application.
\end{frame}

\begin{frame}[fragile,allowframebreaks]{Why secureblue is interesting,
without crowning it}
\protect\phantomsection\label{why-secureblue-is-interesting-without-crowning-it}
secureblue's distinguishing proposition is not that the system is
image-based. Its documented feature set combines allocator hardening
through \texttt{hardened\_malloc}, a specifically confined browser,
removal of legacy privilege paths such as SUID-root programs, restricted
desktop features, and efforts to retain supported sandbox use while
limiting other namespace access (secureblue Project 2026a). That
combination addresses more of the runtime attack surface than
immutability alone, because it hardens the memory-error exploitation
path and the browser --- the two components most exposed to hostile
content --- rather than only the update mechanism. For a user who needs
a recognizable conventional Linux desktop rather than the
compartmentalized workflow of sec.~5, it is the candidate
that most directly engages the exploitation failure path defined in
\end{frame}

\begin{frame}[fragile,allowframebreaks]{Why secureblue is interesting,
without crowning it}
sec.~3.

The release history shows a specific engineering thesis, not a pile of
toggles. Before v4.3.0, the project faced a genuine dilemma it documents
itself: permit unprivileged user namespaces and leave a known
privilege-escalation surface open, or disable them and run Flatpak and
the browser through SUID-root helper binaries --- a privilege path
secureblue was simultaneously trying to eliminate (secureblue Project
2026b). The v4.3.0 resolution moves the decision from file-system bits
into SELinux policy: unconfined processes and the container domain are
denied user namespaces by default, while Flatpak and the browser run in
their own SELinux domains that are granted namespace creation. The
design is in the same family as Ubuntu's restricted unprivileged user
namespaces introduced with Ubuntu 23.10, but stricter in the dimensions
the project chooses to enforce (secureblue Project 2026b). The v4.4.0
\end{frame}

\begin{frame}[fragile,allowframebreaks]{Why secureblue is interesting,
without crowning it}
cycle rebranded the hardened Chromium build as Trivalent --- avoiding
another project's trademark while allowing side-by-side installation
with Fedora's stock Chromium (secureblue Project 2026a) --- and the
v4.9.1 release completed the removal of SUID-root programs, closing the
legacy setuid privilege path the earlier policy work had made redundant.
The direction matters more than any single item: hardening is migrating
from user-configureable exceptions into shipped, policy-enforced
defaults, which is exactly the shape the
\breaktt{application\_confinement} property rewards.

\end{frame}

\begin{frame}[fragile,allowframebreaks]{Why secureblue is interesting,
without crowning it}
The recommendation remains conditional for three documented reasons. The
feature list does not establish independent resistance to a capable
adversary: this review produced no comparative exploit-success or
patch-latency measurements, and project documentation is not a
substitute for either. The compatibility restrictions are real ---
disabled Xwayland and default-deny user namespaces will require
exceptions for some workflows, and each exception is a hand-carved
reduction of the baseline (secureblue Project 2026b). And the strongest
hardening is reversible by the very operator it protects: a hardening
menu that users routinely unwind under compatibility pressure fails the
\texttt{human\_usability} property regardless of its design quality. The
conditional form matters more than the candidate's name: if secureblue's
supported workflow fits the owner's applications, it can be preferable
to assembling a large custom hardening stack; if exceptions proliferate,
\end{frame}

\begin{frame}[fragile,allowframebreaks]{Why secureblue is interesting,
without crowning it}
a more standard environment that is actually maintained may deliver more
security than an abandoned hardened one.
\end{frame}

\begin{frame}[allowframebreaks]{The Flatpak permission-grants warning}
\protect\phantomsection\label{the-flatpak-permission-grants-warning}
Flatpak's documentation describes a restrictive basic sandbox, but also
permission grants that expand access and explicitly warns about
unrestricted bus access on the session or system D-Bus (Flatpak Project
2026). The practical consequence cuts against two common shorthand
judgments. Neither ``uses Flatpak'' nor ``has SELinux'' is evidence that
a particular application is confined to anything the operator would
approve: a permission grant can hand an application access to files,
devices, or the entire message bus, and the grant was usually made once,
at install time, in exchange for the application running at all. The
question that matters --- which files, devices, services, and
credentials this specific application can reach --- is answered by
inspecting the deployed grants, not by the packaging technology's
reputation. This is complete mediation applied to application
\end{frame}

\begin{frame}[allowframebreaks]{The Flatpak permission-grants warning}
installation (Saltzer and Schroeder 1975): every expansion of authority
should pass a checked interface, and the default answer should be the
narrow one. Two developments in this review's candidate set move in that
direction: secureblue's SELinux-scoped domain grants for Flatpak replace
a blanket rule with a per-domain checked interface (secureblue Project
2026b), and Ubuntu 26.04's experimental snap permission prompting moves
the grant decision from install time to use time (Canonical 2026). Both
remain init systems' answers, not agent-era answers: neither mediates
what an authorized agent may do with the access it holds.
\end{frame}

\begin{frame}[fragile,allowframebreaks]{The desktop-versus-agent
authority lesson}
\protect\phantomsection\label{the-desktop-versus-agent-authority-lesson}
Every candidate in tbl.~\ref{tbl:desktops} optimizes the desktop's own
attack surface: how its browser, allocator, update path, and application
sandbox resist exploitation. None of them, as documented, answers the
authorized-misuse path. A shell agent running on even the best-hardened
desktop in this table holds whatever credentials and file access the
operator's daily environment holds; allocator hardening and confined
browsers do not mediate what an autonomous agent with a cloud token may
read, transmit, or change. The lesson generalizes across the whole
candidate set: a hardened desktop is a platform for trust boundaries,
not a substitute for drawing them around autonomous agents. Agent work
with untrusted dependencies belongs in a separate trust domain --- a
disposable VM, a separate container host, or a separate machine --- with
\end{frame}

\begin{frame}[fragile,allowframebreaks]{The desktop-versus-agent
authority lesson}
the authority architecture of sec.~10 governing
what crosses that boundary. The property vocabulary of
sec.~4 makes the distinction exact:
\breaktt{application\_confinement} constrains applications;
\texttt{authority} constrains what a component may already do without
exploiting anything. A desktop can be strong on the first and still
grant an agent catastrophic authority under the second. Scenario-level
guidance for choosing among these candidates, including the conditions
that would change the secureblue-first ordering, appears in
sec.~16.
\end{frame}

\end{document}
